\section{Maintenance Playbook}

Three operations account for nearly all routine upkeep of the data-quality pipeline: adding a new panel round, handling a question that is added to or removed from the questionnaire, and updating the dashboard filters when a skip rule changes. Each is a short, ordered procedure.

\subsection{Adding a New Round}

\begin{procbox}{Add a round (worked example: R9)}
\begin{enumerate}
  \item Bump the round-count constant in four scripts. Add the new round number and label to the \texttt{ROUNDS} list in \rd{build_dq.py}, \rd{build_panel.py}, \rd{build_interviewer.py}, and \rd{build_report.py}. In \rd{build_report.py} also add the round's date label.
  \item In \rd{gen_dashboard.py} and \rd{issue_rollup.py}, change the round range \texttt{range(1, N)} to \texttt{range(1, N+1)}. In \rd{gen_dashboard.py} also extend the JavaScript round arrays, the round labels (\texttt{RLABELS}), the CSV export header and rows, and the Module-00 skip-logic fallback round list.
  \item Register the round's questionnaire. Add the round's Kobo XLSForm to \texttt{KOBO\_FILES} in \rd{parse_kobo.py}, and add the round to \texttt{ROUND\_MAP} in \rd{update_pipeline.py}.
  \item Regenerate everything: \texttt{python3 update\_pipeline.py -{}-all}.
  \item Verify. Confirm the new round appears in the rounds list of \rd{CATI/Analysis/QC/cache/dq_data.json}, and that the round's column and label render in \rd{CATI/Analysis/QC/output/l2phl_dq_dashboard.html}.
\end{enumerate}
\end{procbox}

\begin{notebox}{Why so few edits}
The pooled \texttt{.dta} masters are data-driven by a \texttt{round} column, so the checks themselves need no change: only the round-count constants and the new questionnaire form are touched. A partial, in-progress round shows as provisional. Round~9, for example, held 558 households at hand-over; its completeness flags should be read as ``round not finished,'' not as data problems.
\end{notebox}

\subsection{New or Removed Questions}

\begin{procbox}{Handle a new or removed question}
\begin{enumerate}
  \item Check the round's Kobo XLSForm first. It is the ground truth for which variables exist.
  \item Update the module's authoritative variable list and heatmap rows in \rd{build_dq.py}: add or remove the variable, and set its expected non-missing condition.
  \item In \rd{gen_dashboard.py}, standardize the variable in the module-standardization block so the Tracker, Kobo, and Heatmap panels show it consistently.
  \item Annotate variables that are gated by an upstream preload or skip, so they are not falsely flagged as missing.
  \item Re-run \texttt{python3 update\_pipeline.py -{}-all}.
\end{enumerate}
\end{procbox}

\subsection{Updating Filters When Questionnaire Logic Changes}

\begin{procbox}{Update dashboard filters}
\begin{enumerate}
  \item Add the new round's Kobo XLSForm to \rd{parse_kobo.py}, as in ``Adding a New Round.''
  \item \rd{parse_kobo.py} folds each variable's enclosing repeat or group \texttt{relevant} condition into that variable, so a changed skip rule flows into the dashboard's skip-logic panel automatically once the new form is registered.
  \item Re-run \texttt{python3 update\_pipeline.py -{}-all}. The refreshed skip-logic panel and gated checks reflect the new logic.
\end{enumerate}
\end{procbox}

Because the checks are gated by the Kobo \texttt{relevant} conditions, updating a filter is usually just registering the new form, not editing check code by hand.
